\begin{abstract}
本文提出了一种多相机3D目标检测（Multi-Camera 3D Object Detection）框架。
与直接从单目图像估计3D边界框或使用深度预测网络从2D信息生成3D目标检测输入的现有方法不同，本文方法直接在3D空间中进行预测。本文架构从多个相机图像中提取2D特征，然后使用稀疏的3D目标查询集合索引这些2D特征，通过相机变换矩阵将3D位置与多视角图像关联起来。最终，本文模型对每个目标查询进行边界框预测，并使用集合到集合（Set-to-Set）损失来衡量预测与真值之间的差异。
本文研究了连接2D观测与3D预测的两种方案：一种是依赖于密集深度预测和3D场景结构恢复的自底向上（Bottom-Up）方法；另一种是从稀疏目标先验出发，通过反向投影（Backward Projection）查询2D图像特征的自顶向下（Top-Down）方法。
这种自顶向下方法优于自底向上方法（后者在逐像素深度估计后进行目标边界框预测），因为它不受深度预测模型引入的累积误差影响。此外，本文方法不需要非极大值抑制（Non-Maximum Suppression, NMS）等后处理，显著提升了推理速度。本文在nuScenes自动驾驶基准上取得了当前最优（State-of-the-Art, SOTA）的性能。
\end{abstract}

